\documentclass[11pt]{article}
\usepackage[a4paper,margin=27mm]{geometry}
\usepackage{amsmath,amssymb,amsthm}
\usepackage[T1]{fontenc}
\usepackage{lmodern}
\usepackage[hidelinks]{hyperref}
\newtheorem{theorem}{Theorem}
\newtheorem{lemma}{Lemma}
\setlength{\emergencystretch}{3em}
\newcommand{\R}{\mathbb R}
\newcommand{\N}{\mathbb N}
\newcommand{\cl}{\operatorname{cl}}
\title{A numerical obstruction to Conjecture 00000007914}
\author{C0ldSmi1e}
\date{9 October 2026}
\begin{document}
\maketitle
\begin{abstract}
The asserted asymptotic for the increasing finite ranks of a covolume
spectrum is impossible. We prove a stronger statement for every set of
nonnegative real numbers, allowing zero and every real leading constant.
The contradiction concerns the numerical assertions themselves and does
not use any classification of arithmetic Fuchsian or triangle groups.
The accompanying Lean 4 project checks the exact rank argument, the
logarithmic asymptotic, and the nonnegativity interfaces for volumes and limits.
\end{abstract}

\section{Original assertion and scope}
The English original defines $S_{\rm arith}$ as the ``limit-from-below
spectrum'' of covolumes $\operatorname{vol}(\Gamma\backslash H)$ of
arithmetic Fuchsian groups. It asserts that $S_{\rm arith}$ is well ordered
(no descending chains), and that its $k$-th element satisfies
\begin{equation}\label{eq:source}
 v_k\sim c\,k^{-1}(\log k)^{-2/3}\qquad(k\longrightarrow\infty).
\end{equation}
It additionally asserts that the smallest noncompact arithmetic triangle
group is $(2,3,7)$ and makes a next-smallest trichotomy claim with covolume
ratio $3/7$. Both language versions contain the same displayed asymptotic
and ordered-spectrum assertions. The exact bilingual source is included
as \texttt{conjecture.md}.

This report disproves the conjunction by refuting its finite-rank
asymptotic. It makes no separate claim about the triangle groups,
the trichotomy, or well-ordering in isolation. Throughout, an element's
rank means its rank in the inherited increasing numerical order, as in
the original well-order assertion. The formula is treated as an
asymptotic for the elements themselves, not for a reciprocal or a
counting function.

\section{Universal obstruction}
For $S\subseteq\R$, say that $x$ has finite rank $n$ if $x\in S$ and the
\emph{entire} predecessor set
\[
 P_S(x)=\{y\in S:y<x\}
\]
is finite and has cardinality $n$. The first element has rank zero.
This definition neither counts groups with repeated equal covolumes nor
requires all of $S$ to have finite rank.

\begin{lemma}\label{lem:rank}
If $v_k$ has rank $k-1$ in $S$ for each integer $k\geq1$, then
$v_m<v_n$ whenever $1\leq m<n$.
\end{lemma}
\begin{proof}
If $v_n\leq v_m$, then $P_S(v_n)\subseteq P_S(v_m)$.
Their finite cardinalities would give $n-1\leq m-1$, a contradiction.
\end{proof}

For a fixed $c\in\R$, put
\[
 g_c(k)=c\,k^{-1}(\log k)^{-2/3}\quad(k\geq2).
\]
The logarithm is natural and the power is the usual real power of the
positive number $\log k$.
\begin{lemma}\label{lem:limit}
For every $c\in\R$, $g_c(k)\longrightarrow0$ as $k\longrightarrow\infty$.
If $c\ne0$ and $v_k/g_c(k)\longrightarrow1$, then $v_k\longrightarrow0$.
\end{lemma}
\begin{proof}
The reciprocal $k^{-1}$ tends to zero. Also $\log k$ tends to positive
infinity, so $(\log k)^{-2/3}$ tends to zero. The product, multiplied by
the fixed constant $c$, therefore tends to zero. In the second assertion,
write $v_k=(v_k/g_c(k))g_c(k)$ for $k\geq2$ and take limits.
\end{proof}

\begin{theorem}\label{thm:main}
Let $S\subseteq[0,\infty)$. There is no sequence whose $k$-th term has
rank $k-1$ in $S$ for every $k\geq1$ and which satisfies
\eqref{eq:source} in the ordinary ratio-limit sense, for any nonzero
real constant $c$.
\end{theorem}
\begin{proof}
Lemma~\ref{lem:rank} gives $0\leq v_1<v_2$ and $v_k\geq v_2>0$
for every $k\geq2$. Lemma~\ref{lem:limit} gives $v_k\to0$.
Passing to the limit in the eventual inequality $v_2\leq v_k$ gives
$v_2\leq0$, contradicting $v_2>0$.
\end{proof}

Zero is allowed in $S$: it is strict separation of the first two finite
ranks that supplies a positive lower bound. No monotonicity of the
comparison function is transferred through asymptotic equivalence; only
its limit is transferred. No exhaustion of a transfinite well-order by a
sequence is assumed. The asserted asymptotic presupposes arbitrarily
large finite ranks, which is all that is relevant here.

\section{Why the obstruction applies to covolumes}
A covolume is a nonnegative measure value. Finite covolumes can therefore
be regarded as nonnegative real numbers without any geometric
classification or positive lower-bound theorem. More explicitly, let
$(X_i,\mu_i)$ be \emph{any} family of measured spaces and put
\[
 C=\{\mu_i(X_i):\mu_i(X_i)<\infty\}\subseteq\R.
\]
Then $C\subseteq[0,\infty)$. The Lean interface uses extended-nonnegative
real measures and excludes the infinite value before conversion to a
real number. In particular, it never turns an infinite covolume into a
spurious real zero.

The source does not specify the operation called ``limit-from-below
spectrum.'' We do not choose one favorable construction. The result
holds for every nonnegative real set, and in particular for every
$S\subseteq\cl(C)$: the closed half-line $[0,\infty)$ contains $C$ and
therefore its closure. This covers raw covolumes, their closure,
one-sided accumulation values, and arbitrary selections of such values.
Strict positivity of a raw covolume is \emph{not} promoted to strict
positivity of a limit.

The Chinese wording can also be read as a lower-limit (\emph{liminf})
spectrum. A finite liminf of a sequence of nonnegative covolumes remains
nonnegative: every tail has nonnegative infimum, so the supremum of
those tail infima is nonnegative. Thus the same obstruction applies to
that reading as well. Infinite spectrum values cannot serve as terms of
the stated real asymptotic. No interpretation that retains covolumes or
their ordinary finite limits can introduce negative real values.

Consequently the actual arithmetic covolumes and the ordinary limiting
readings of the source all satisfy the necessary nonnegative condition.
Theorem~\ref{thm:main} rules out the displayed rank asymptotic for any such
spectrum. Conjoining further well-ordering or triangle-group assertions
cannot repair the contradiction.

\clearpage
\section{Formal statement and verification interface}
The project pins Lean 4.19.0 and Mathlib v4.19.0 at revision
\begin{center}\small\texttt{c44e0c8ee63ca166450922a373c7409c5d26b00b}\end{center}
The main file is \texttt{lean/CovolumeSpectrum.lean}.

\begin{itemize}
\item \texttt{HasRank} uses a finite set whose membership is exactly
membership in the full predecessor set, and whose cardinality is the
specified rank. \texttt{rank\_strictMono} proves Lemma~\ref{lem:rank}.
\item \texttt{scale} is the literal displayed reciprocal-logarithmic
expression with exponent $-(2/3)$ in $\R$.
\texttt{scale\_tendsto\_zero} proves its limit.
\item \texttt{asymptotic\_iff\_ratio} identifies Mathlib's asymptotic
equivalence with the ordinary ratio limit whenever $c\ne0$.
The totalized values at indices zero and one are irrelevant to the
filter at infinity.
\item \texttt{no\_rankedAsymptotic} proves the universal obstruction.
It even allows every $c\in\R$ under Mathlib's extended definition
$v-g=o(g)$. For $c=0$, this requires eventual equality to zero and is
also impossible; the ordinary ratio definition itself excludes $c=0$.
\item \texttt{no\_numericalConjecture} also includes the original
well-foundedness requirement. \texttt{no\_full\_conjunction} permits
an arbitrary proposition for all remaining source clauses.
\item \texttt{finiteMeasureValues\_nonnegative}\\
\texttt{no\_measure\_spectrum\_conjunction}\\
These lemmas supply the general
measure-volume interface and the closure-of-volume-values conclusion.
The geometric quotient construction is not assumed as an additional
theorem: the universal obstruction applies to any family of measures.
\item \texttt{finite\_liminf\_nonnegative} works with the
extended-real liminf and an explicitly finite real value.
\texttt{no\_liminf\_spectrum\_conjunction} handles selections of finite
sequential lower limits directly, without imposing closure as an extra
assumption or requiring a separate boundedness hypothesis.
\end{itemize}

The source-to-formal correspondence uses only the ordinary meanings of
volume, finite numerical rank, and asymptotic equivalence. The
unnecessary arithmetic classifications are not formalized and are not
claimed as verified. The exact commands, theorem-dependency output,
independent review, and final validation records accompany the package.
These are local verification results; maintainer acceptance is separate.

\end{document}
